\documentclass[11pt,a4paper]{article}
\usepackage[T1]{fontenc}
\usepackage[utf8]{inputenc}
\usepackage{lmodern,amsmath,amssymb}
\usepackage[margin=25mm]{geometry}
\usepackage[hidelinks]{hyperref}
\hypersetup{pdftitle={Size stabilization and action normalization},pdfauthor={navstokgap research working notes}}
\newcommand{\tightlist}{\setlength{\itemsep}{0pt}\setlength{\parskip}{0pt}}
\setlength{\emergencystretch}{3em}
\setcounter{secnumdepth}{0}
\begin{document}
\hypertarget{stabilizing-a-topological-radius-exposes-an-action-coefficient}{%
\section{Stabilizing a topological radius exposes an action
coefficient}\label{stabilizing-a-topological-radius-exposes-an-action-coefficient}}

A quadratic/quartic field energy has a finite preferred scale within an
explicit degree-one family. Its energy-times-radius divided by wave
speed is proportional to the quartic coefficient divided by that speed.
Scaling both field coefficients leaves the classical dynamics unchanged
while scaling this action. Size stabilization succeeds; absolute action
normalization remains an independent input.

Status: Q05 exploratory derivation, not an accepted-ledger claim.
Independent review and a bounded librarian comparison precede any
promotion. The scale minimum below is a trial-family result, not a proof
of a full field minimizer.

\hypertarget{field-sector-and-kinetic-term}{%
\subsection{1. Field, sector and kinetic
term}\label{field-sector-and-kinetic-term}}

Let \(n:\mathbb R^3\to S^3\subset\mathbb R^4\) be smooth, approach a
fixed point at spatial infinity, and have finite energy

\[E[n]=\int\left[a\sum_i|\partial_i n|^2+
b\sum_{i<j}\left(|\partial_i n|^2|\partial_j n|^2
-(\partial_i n\cdot\partial_j n)^2\right)\right]d^3x. \tag{1}\]

Here \(a>0\) has units energy/length, \(b>0\) has units energy times
length, and \(c>0\) has speed units. Compactifying the domain defines
the integer degree \(N\). This is the massless Skyrme-type
quadratic/quartic construction written in real sphere coordinates; its
standard ingredients have prior art.

To specify physical time, choose the conservative Lagrangian

\[L=\int\left[\frac{a}{c^2}|\dot n|^2+
\frac{b}{c^2}\sum_i\left(|\dot n|^2|\partial_i n|^2
-(\dot n\cdot\partial_i n)^2\right)\right]d^3x-E[n]. \tag{2}\]

The velocity quadratic form is positive. This is the usual
Lorentz-invariant quadratic/quartic completion, with coordinate
\(x^0=ct\); all equations here use physical time \(t\). We require
smooth solutions where they exist. This specification supplies neither
global nonlinear well-posedness nor a hard speed theorem for arbitrary
backgrounds.

For completeness, if \(\lambda_i\ge0\) are the singular values of
\(dn\), the static density is
\(a\sum_i\lambda_i^2+b\sum_{i<j}\lambda_i^2\lambda_j^2\). Pairing
\(a\lambda_i^2\) with \(b\lambda_j^2\lambda_k^2\) and applying the
arithmetic-geometric mean inequality gives density at least
\(6\sqrt{ab}\lambda_1\lambda_2\lambda_3\). Since the unit three-sphere
has volume \(2\pi^2\), integration yields the standard degree estimate

\[E\ge12\pi^2\sqrt{ab}\,|N|. \tag{3}\]

The absolute Jacobian integral bounds the absolute degree even when
local orientations vary. A positive sector energy is distinct from a
vacuum excitation gap or an action bound.

\hypertarget{an-explicit-size-test}{%
\subsection{2. An explicit size test}\label{an-explicit-size-test}}

Write \(y=x/R\), \(r=|y|\), and use the inverse stereographic profile

\[n_R(x)=\frac{(2y_1,2y_2,2y_3,r^2-1)}{1+r^2},\qquad R>0.\]

It has \(|N|=1\), and \(R\) is the sphere's equator radius in the
domain. All three dimensionless singular values are \(2/(1+r^2)\). Thus

\[A=48\pi\int_0^\infty\frac{r^2\,dr}{(1+r^2)^2}=12\pi^2,
\qquad
B=192\pi\int_0^\infty\frac{r^2\,dr}{(1+r^2)^4}=6\pi^2.\]

The integrals follow directly from \(r=\tan u\), \(0<u<\pi/2\). Spatial
scaling gives

\[E(R)=aAR+\frac{bB}{R},\qquad
R_* =\sqrt{\frac{bB}{aA}}=\sqrt{\frac{b}{2a}},\qquad
E(R_*)=12\sqrt2\pi^2\sqrt{ab}. \tag{4}\]

Unlike Q04, shrinking this family sends its energy to infinity.
Expansion also costs energy. This proves a unique minimum along its
scale direction, with positive second derivative \(2bB/R^3\). General
shape perturbations remain outside this one-parameter test. In
particular, \(n_{R_*}\) is a trial field, not asserted to solve the
static Euler--Lagrange equation.

\hypertarget{where-the-action-enters}{%
\subsection{3. Where the action enters}\label{where-the-action-enters}}

Define the same type of diagnostic as Q04: evaluate the magnitude of
static Lagrangian action over the profile's crossing duration
\(T_R=R/c\),

\[\mathcal A_R=E(R)R/c=(aAR^2+bB)/c. \tag{5}\]

For off-shell static trial fields this is a functional diagnostic, not
an on-shell trajectory action. It has infimum \(6\pi^2 b/c\) within this
family; at its energy-optimal scale it equals

\[\mathcal A_{R_*}=12\pi^2 b/c. \tag{6}\]

The product \(b/c\) already has action units. More generally, setting
\(x=\ell y\), \(t=\ell s/c\) with \(\ell=\sqrt{b/a}\) writes the full
action as

\[S=\frac{b}{c}\,\widetilde S[n], \tag{7}\]

where the dimensionless functional has unit quadratic and quartic
coefficients. This includes the time-derivative terms, not just static
energy. It identifies exactly which normalization a proposed quantum
phase would use.

There is also an exact countertest. Replace \((a,b)\) by
\((\eta a,\eta b)\), \(\eta>0\). The length \(\ell\), field equations,
degree, and characteristic vacuum speed remain unchanged. Energies and
actions multiply by \(\eta\). Every positive \(\eta\) defines a
nondegenerate model, and actions tend to zero as \(\eta\downarrow0\). No
claim about dynamics at the degenerate endpoint is needed. This varies
the physical stiffnesses; it does not deny fixed-coefficient sector
bounds. It shows that field equations and topology alone do not fix
their common mechanical normalization.

\hypertarget{physical-spectrum-and-next-decision}{%
\subsection{4. Physical spectrum and next
decision}\label{physical-spectrum-and-next-decision}}

Near a constant vacuum, write \(n=(\pi,\sqrt{1-|\pi|^2})\). To quadratic
order,

\[L^{(2)}=a\int[c^{-2}|\dot\pi|^2-|\nabla\pi|^2]d^3x,
\qquad \Omega(k)=c|k|.\]

The quartic term has no quadratic vacuum contribution. On infinite space
there are arbitrarily low frequencies. Fixing a nonzero topological
sector and stabilizing its radius therefore leaves this vacuum spectrum
gapless. It also supplies no mechanism forcing a trajectory to carry
nonzero degree.

Retain gradient competition as a conservative size-exclusion mechanism.
Its action coefficient is explicit, while the vacuum-gap question
remains separate. After two topology steps, park further adjustable
stabilizers. The next physical test should couple sectors that initially
have independent normalizations: determine whether reciprocal energy
exchange can constrain their ratio and whether one overall action
normalization survives. This can turn the common-normalization
obstruction into a concrete composition test, rather than adding another
term whose coefficient contains the answer.

\hypertarget{source-boundary}{%
\subsection{Source boundary}\label{source-boundary}}

The discovery lead is \emph{Topological energy bounds in generalized
Skyrme models}, \href{https://arxiv.org/abs/1311.2939}{arXiv:1311.2939},
\href{https://journals.aps.org/prd/abstract/10.1103/PhysRevD.89.065010}{publisher
abstract}. Coverage here is the publisher/search abstract only: it
identifies opposite Derrick scaling and topological bounds as
established ingredients. Full-text formula matching and independent
review are pending, not silently imported. The normalization, trial
integrals, kinetic expansion and action comparison above are written
calculations. No computational verification scripts were used.
\end{document}
